
\begin{exercise}[Long term investment returns \Coffeecup]
    In this exercise we restrict ourselves to a single asset.
    That is \(n=1\) and we write \(R_t = R_t^1\) for the return of the asset at
    time \(t\), cf.~\eqref{eq: return definition}.
    Assume that the returns \(R_0, R_1, \dots\) are independent and identically
    distributed with \(\E{|\ln(R_1)|} < \infty\). Compare the
    geometric mean with the arithmetic mean, that is
    \[
        \Bigl(\prod_{t=0}^{n-1} R_t\Bigr)^\frac1n
        \quad \text{vs.} \quad
        \frac1n\sum_{t=0}^{n-1} R_t.
    \]
    Which one captures the long term growth of the investment better?
    \begin{solution}
        The geometric mean of the returns \(\mean{R}^g_n \coloneq \Bigl(\prod_{t=0}^{n-1} R_t\Bigr)^\frac1n\)
        repeated for \(n\) periods gives the same return as the actual return over \(n\) periods:
        \[
            (\mean{R}^g_n)^n = \prod_{t=0}^{n-1} R_t = \frac{P_n}{P_0}.
        \]
        Consequently, \(r^g \coloneq \exp(\E{\ln(R_1)})\) is the long term
        growth rate that stems from the law of large numbers
        \[
            \mean{R}^g_n  = \exp\Bigl(\frac1n \sum_{t=0}^{n-1} \ln(R_t)\Bigr) \overset{\as}\to r^g.
        \]
        In contrast, \(\E{R_1}\) may be arbitrarily far from \(r^g\) and
        \(\E{R_1}^n\) is thereby not a good approximation of the return over \(n\) periods.
        The arithmetic mean approximates \(\E{R_1}\) and is thereby not a good
        approximation of the long term growth of the investment.
    \end{solution}
\end{exercise}